%---------------------------------------------------------%
%   CHAPTER 3: PSM REPORT CONTENT OUTLINE                  %
%---------------------------------------------------------%
\chapter{PSM Report Content Outline}

This chapter adapts the PSM content guideline \citep{psmguideline} into a clean, numbered outline. It follows the chapter structure of the UMP Thesis Guide \citep{umpthesis2019}.

\section{Report Structure Overview}

The PSM 1 report (proposal stage) normally contains three chapters plus references:

\begin{enumerate}
    \item Chapter 1: Introduction
    \item Chapter 2: Literature Review
    \item Chapter 3: Methodology
\end{enumerate}

The final PSM report extends to the five-chapter structure recommended by the UMP Thesis Guide: Introduction, Literature Review, Research Design and Methodology, Results and Discussion, and Conclusion and Recommendations. Students should confirm the exact chapter division with their supervisor.

\section{Chapter 1: Introduction}

\subsection{Introduction}

Describe the project briefly and explain the related issues. Ease the reader into the topic: state the general topic, give background, review related literature briefly, define the terms and scope, and outline the current situation.

\subsection{Problem Statement}

State the current issue that the project seeks to address. A good problem statement answers \emph{who, what, when, where}, and \emph{why}. Recommended elements:

\begin{itemize}
    \item Describe how things should work.
    \item Explain the problem and why it matters.
    \item Support the claim with cited sources.
    \item Propose a solution.
    \item Explain the benefits of the proposed solution.
    \item Conclude by summarising the problem and the solution.
\end{itemize}

\subsection{Objective}

Objectives are specific actions to be achieved, aligned with the project goal. They should follow the SMART concept: Specific, Measurable, Achievable, Realistic, and Time-framed. Write objectives as a list using roman numbering (i, ii, iii) and active statements. Example:

\begin{verbatim}
The objectives of this project are:
(i) To study ...
(ii) To develop ...
(iii) To evaluate ...
\end{verbatim}

\subsection{Scope}

Explain the boundary of the project and its users: the extent of the work and the parameters within which the study operates. Write the scope in list form (i, ii, iii) and order the items according to the objectives.

\subsection{Report Organization}

Summarise each chapter of the report. Example:

\begin{verbatim}
This report consists of five chapters. Chapter 1 discusses the
introduction to the project. Chapter 2 explains and discusses ...
\end{verbatim}

\section{Chapter 2: Literature Review}

\subsection{Introduction}

Explain the existing problem or solution from previous work. Describe in detail the technique, method, hardware, or technology suitable for the project. The review should answer \emph{What, Why,} and \emph{How}, and comparisons may be shown in tables. State the sources clearly (book, journal, conference, proceedings, thesis, or website). Citations and bibliography follow the UMP Thesis Guide APA style.

\subsection{Existing Systems or Works}

Describe at least three existing systems or works that are relevant to the proposed solution:

\begin{enumerate}
    \item \textbf{System A} --- explanation with supporting diagram or interface.
    \item \textbf{System B} --- explanation with supporting diagram or interface.
    \item \textbf{System C} --- explanation with supporting diagram or interface.
\end{enumerate}

\subsection{Analysis and Comparison of Existing Systems}

Compare the three existing systems and review the comparison in terms of advantages and disadvantages. Then discuss how relevant, significant, or impactful the comparison is to the project title.

\subsection{Summary}

Summarise the key findings of Chapter 2.

\section{Chapter 3: Methodology}

\subsection{Introduction}

Discuss the overall approach or framework of the undergraduate project. The chapter should cover the method, technique, or approach used, and discuss the methodology in detail so the project can be accomplished.

\subsection{Project Management Framework or Methodology}

Discuss and justify the selected development or research methodology and its phases. Project-based work may use SDLC (Agile, RAD, etc.); research-based work uses a research framework. Describe the activities in each phase and how they relate to the project.

\subsection{Project Requirements}

State what is needed to complete the project.

\begin{itemize}
    \item \textbf{Project-based:} functional and non-functional requirements, constraints, and limitations. User requirements may come from interviews, surveys, client feedback, or observation of the existing manual process.
    \item \textbf{Research-based:} input, output, process description, constraints and limitations, and case study.
\end{itemize}

\subsection{Proposed Design}

Present the proposed design based on the requirements.

\begin{itemize}
    \item \textbf{Project-based:} flowchart, context diagram, use case diagram and description, activity diagram, storyboard.
    \item \textbf{Research-based:} pseudocode, algorithm, flowchart, model, or a related design depending on the type of research.
\end{itemize}

\subsection{Data Design}

Describe the data or assets involved in the project.

\begin{itemize}
    \item \textbf{Project-based:} ERD, database dictionary (primary key, foreign key), assets, and the process of creating the database.
    \item \textbf{Research-based:} dataset description and justification of the selected dataset.
\end{itemize}

\subsection{Proof of Initial Concept}

For PSM 1, provide at least a proof of concept (POC) that verifies the idea can be executed from a technical point of view.

\begin{itemize}
    \item \textbf{Project-based:} design of the proof of concept with explanation.
    \item \textbf{Research-based:} evidence of early work, experiment, or preliminary findings.
\end{itemize}

\subsection{Testing or Validation Plan}

Plan the testing or validation process clearly. Examples include selecting data for testing, how testing is implemented, and the measurement or evaluation used (e.g., user acceptance testing, comparison testing). Project-based work describes the testing strategy for the proposed solution; research-based work covers a simple set of inputs, processes, outputs, and comparison of results.

\subsection{Potential Use of Proposed Solution}

Explain the potential use of the proposed solution in real situations or its potential for commercialisation. Discuss the benefits and significance of the project in daily life.

\subsection{Gantt Chart}

Show the project phases from start to completion. For PSM 1, show the estimated duration. The chart should follow the selected development methodology activities, not the report chapters.

\section{References}

References may come from books, journals, conference proceedings, magazines, or websites. Record all references to identify the copyright and originality of the work. Format citations and bibliography according to the UMP Thesis Guide APA style \citep{umpthesis2019}.

\section{Notes for PSM 2 (Final Report)}

For the final report, add:

\begin{itemize}
    \item \textbf{Chapter 4: Results and Discussion} --- present and interpret the results, verify that all figures and tables are necessary and appropriate, and discuss the findings in the context of previous work.
    \item \textbf{Chapter 5: Conclusion and Recommendations} --- summarise the findings, state the contributions and limitations, and propose future work.
\end{itemize}
